\documentclass[12pt]{article}
\usepackage[utf8]{inputenc}
\usepackage[T1]{fontenc}
\usepackage[margin=1in]{geometry}
\usepackage{setspace}
\usepackage{xcolor}
\usepackage[colorlinks=true, urlcolor=blue!60!black, linkcolor=blue!60!black]{hyperref}
\setlength{\parindent}{0em}
\setlength{\parskip}{0.45em}
\renewcommand{\baselinestretch}{1.0}
\begin{document}
\begin{center}
{\Large \textit{U.S. Term Limits, Inc. v. Thornton}, 514 U.S. 779 (1995)\par}
\vspace{0.3em}
{\normalsize Case Brief --- One-Page Edition\par}
\end{center}

\textbf{Memory Jogger.} No state-imposed congressional term limits, even as ballot-access rules; only Article V can add qualifications (\href{https://supreme.justia.com/cases/federal/us/514/779/}{514 U.S. 779}).

\textbf{Facts.} Arkansas's 1992 Amendment 73 kept anyone elected to three House terms or two Senate terms off the general-election ballot, leaving only write-in campaigns, to combat ``[e]ntrenched incumbency.''

\textbf{Procedural History.} Voter Bobbie Hill and the League of Women Voters sued in state court, joined by Representative Ray Thornton; the circuit court and a 5--2 Arkansas Supreme Court held \S~3 unconstitutional. Twenty-two States had similar measures.

\textbf{Votes.} Affirmed 5--4. Stevens wrote for the Court; Kennedy concurred; Thomas, with Rehnquist, O'Connor, and Scalia, dissented.

\textbf{Holding.} The Qualifications Clauses are fixed and exclusive; States have no reserved power to add qualifications and may not use the Elections Clause to handicap a class of candidates.

\textbf{Reasoning.} The Tenth Amendment reserves only pre-existing powers, and Congress did not exist before 1787; members are paid and judged nationally, and the Framers feared state interference. \S~3 had the ``sole purpose of creating additional qualifications indirectly.'' Kennedy stressed dual sovereignty; Thomas argued the Constitution rests on the people of each State, the Clauses set only floors, and where the Constitution is silent the States retain power.

\textbf{Examples.} \emph{Same-Side:} a district-residency rule for House candidates. \emph{Opposite-Side:} term limits on state officials. \emph{Fence-Sitter:} a ballot notation of terms served.

\textbf{Critique.} The case turns on whether one national people or the separate peoples of the States made the Constitution; meanwhile States cannot offset incumbents' self-granted advantages, and Congress has little incentive to limit itself.

\textbf{Key Quotation.} ``Federalism was our Nation's own discovery. The Framers split the atom of sovereignty. It was the genius of their idea that our citizens would have two political capacities, one state and one federal, each protected from incursion by the other.'' --- Justice Kennedy, concurring, 514 U.S. at 838.

\end{document}
